% Part: computability
% Chapter: recursive-functions
% Section: examples

\documentclass[../../../include/open-logic-section]{subfiles}

\begin{document}

\olfileid{cmp}{rec}{exa}
\olsection{原始递归函数的例子}

我们已经知道一些原始递归函数的例子：加法函数~$\Add$ 和乘法函数~$\Mult$。
恒等函数 $\fn{id}(x) = x$ 是原始递归的，因为它就是~$\Proj{1}{0}$。
常数函数 $\fn{const}_n(x) = n$ 也是原始递归的，因为它们可以由 $\Zero$ 和 $\Succ$ 通过反复复合定义出来。
当我们在原始递归定义中需要使用常数时，这会很有用。例如，若要定义函数 $f(x) = 2 \cdot x$，可以通过将常数函数 $\fn{const}_n(x)$ 与乘法函数复合来得到，即 $f(x) =
\Mult(\fn{const}_2(x), \Proj{1}{0}(x))$。
今后我们将经常使用这一技巧。

\begin{prop}
  指数函数 $\fn{exp}(x, y) = x^y$ 是原始递归的。
\end{prop}

\begin{proof}
  我们可以将 $\fn{exp}$ 原始递归地定义为
  \begin{align*}
    \fn{exp}(x, 0) & = 1\\
    \fn{exp}(x, y+1) & = \Mult(x, \fn{exp}(x,y)).
    \intertext{严格来说，这并不是一个由原始递归函数出发的递归定义。
      不过，严格按定义来说，我们有：}
    \fn{exp}(x, 0) & = f(x)\\
    \fn{exp}(x, y+1) & = g(x, y, \fn{exp}(x,y)).
    \intertext{其中}
    f(x) & = \Succ(\Zero(x)) = 1\\
    g(x, y, z) & = \Mult(\Proj{3}{0}(x, y, z), \Proj{3}{2}(x, y, z)) = x \cdot z
  \end{align*}
  因此 $f$ 和 $g$ 是由原始递归函数通过复合定义的。
\end{proof}

\begin{prop}
  由
  \[
  \fn{pred}(y) = \begin{cases}
    0 & \text{若 $y=0$}\\
    y-1 & \text{否则}
  \end{cases}
  \]
  定义的前驱函数 $\fn{pred}(y)$ 是原始递归的。
\end{prop}

\begin{proof}
  注意到
  \begin{align*}
   \fn{pred}(0) & = 0 \text{ 且}\\
   \fn{pred}(y+1) & = y.
 \end{align*}
  这几乎就是一个原始递归定义。严格地说，它并不完全符合原始递归的定义模式，因为该模式要求至少有一个额外的自变元~$x$。此外，它也有些反常，因为在定义 $\fn{pred}(y+1)$ 时实际上并没有用到 $\fn{pred}(y)$。但我们可以先通过
  \begin{align*}
   \fn{pred}'(x, 0) & = \Zero(x) = 0,\\
   \fn{pred}'(x, y+1) & = \Proj{3}{1}(x, y, \fn{pred'}(x, y)) = y.
 \end{align*}
定义 $\fn{pred}'(x, y)$，
然后再通过复合由它定义 $\fn{pred}$，例如，定义为
$\fn{pred}(x) = \fn{pred}'(\Zero(x), \Proj{1}{0}(x))$。
\end{proof}

\begin{prop}
  阶乘函数 $\fn{fac}(x) = \fact{x} = 1 \cdot 2 \cdot 3
  \cdot \dots \cdot x$ 是原始递归的。
\end{prop}

\begin{proof}
  最直接的原始递归定义是
  \begin{align*}
    \fn{fac}(0) &= 1\\
    \fn{fac}(y+1) & = \fn{fac}(y) \cdot (y+1).
    \intertext{严格形式地，我们必须先定义一个二元函数 $h$，}
    h(x, 0) & = \fn{const}_1(x)\\
    h(x, y+1) & = g(x, y, h(x, y))
    \intertext{其中 $g(x, y, z) = \Mult(\Proj{3}{2}(x, y, z),
      \Succ(\Proj{3}{1}(x, y, z)))$，然后令}
    \fn{fac}(y) & = h(\Proj{1}{0}(y), \Proj{1}{0}(y)) = h(y,y).
  \end{align*}
  从现在开始我们会更随意一些，不再详细给出通过复合和原始递归的严格形式定义。
\end{proof}

\begin{prop}
  由
  \[
  x \tsub y = \begin{cases}
    0 & \text{若 $x < y$}\\
    x-y & \text{否则}
  \end{cases}
  \]
  定义的截断减法 $x \tsub y$ 是原始递归的。
\end{prop}

\begin{proof}
  我们有：
  \begin{align*}
    x \tsub 0 & = x\\
    x \tsub (y+1) & = \fn{pred}(x \tsub y) 
  \end{align*}
\end{proof}

\begin{prop}
 $x$ 与 $y$ 之间的距离 $\left|x-y\right|$ 是原始递归的。
\end{prop}

\begin{proof}
  我们有 $\left| x-y \right| = (x \tsub y) + (y \tsub x)$，因此距离函数可以由 $+$ 和 $\tsub$ 复合定义得到，而 $+$ 和 $\tsub$ 都是原始递归的。
\end{proof}

\begin{prop}
 $x$ 与 $y$ 的最大值 $\fn{max}(x,y)$ 是原始递归的。
\end{prop}

\begin{proof}
  我们可以通过将 $+$ 和 $\tsub$ 复合来定义 $\fn{max}(x,y)$：
  \[
  \fn{max}(x,y) \defis x + (y \tsub x).
  \]
  若 $x$ 是最大值，即 $x \ge y$，则 $y \tsub x = 0$，因此 $x
  + (y \tsub x) = x + 0 = x$。若 $y$ 是最大值，则 $y \tsub x = y - x$，于是 $x + (y \tsub x) = x + (y - x) = y$。
\end{proof}

\begin{prop}
\ollabel{prop:min-pr}
 $x$ 与 $y$ 的最小值 $\fn{min}(x,y)$ 是原始递归的。
\end{prop}

\begin{proof}
  练习。
\end{proof}

\begin{prob}
证明 \olref[cmp][rec][exa]{prop:min-pr}。
\end{prob}

\begin{prob}
证明 \[
f(x, y) =
2^{(2^{\iddots^{2^{x}}})}\raisebox{1ex}{\bigg\rbrace}
\raisebox{1ex}{\text {$y$ $2$'s}}\] 是原始递归的。
\end{prob}

\begin{prob}
证明整数除法 $d(x, y) = \lfloor x/y \rfloor$（即舍去小数点后所有部分的除法）是原始递归的。当 $y = 0$ 时，我们规定 $d(x, y) = 0$。请使用原始递归与复合给出 $d$ 的一个显式定义。
\end{prob}

\begin{prop}
原始递归函数集在以下两种运算下封闭：
\begin{enumerate}
\item 有限和：若 $f(\vec x, z)$ 是原始递归的，则函数
\[
g(\vec x, y) \defis \sum_{z = 0}^y f(\vec x, z).
\]
也是原始递归的。
\item 有限积：若 $f(\vec x, z)$ 是原始递归的，则函数
\[
h(\vec x, y) \defis \prod_{z = 0}^y f(\vec x, z).
\]
也是原始递归的。
\end{enumerate}
\end{prop}

\begin{proof}
例如，有限和可通过下列递归方程定义：
\begin{align*}
  g(\vec x, 0) & = f(\vec x, 0)\\
  g(\vec x, y+1) & = g(\vec x, y) + f(\vec x, y+1).
\end{align*}
\end{proof}


\end{document}